\chapter{刚体转动}

\section{平移之外，还要跟踪取向}
刚体近似要求各质点间距离不随时间变化，适用于形变对所研究运动可忽略的情形。
把参考点选在质心，任一点的速度为
\begin{equation}
 \bm v_a=\dot{\bm R}+\bm\omega\times\bm\rho_a.
\end{equation}
利用质心关系消去交叉项，动能分解为
\begin{equation}
 T=\frac12M\dot{\bm R}^{\,2}+\frac12\bm\omega\tp I\bm\omega,\qquad
 I_{ij}=\sum_a m_a(\rho_a^2\delta_{ij}-\rho_{ai}\rho_{aj}).
 \label{eq:cm-inertia}
\end{equation}
连续分布时将求和换为$\int\d m$。惯性张量是实对称矩阵，
其方向依赖性描述在同一角速率下绕不同方向转动所需的动能。
关于同一参考点，$\bm\ell=I\bm\omega$，角动量一般不与角速度平行。

主轴是$I$的正交特征矢量，主惯量为$I_1,I_2,I_3$。
对通常三维刚体它们为正；理想直线质量分布绕自身轴的惯量可以为零。
关于与质心相距$\bm a$的点、保持平行坐标轴，平行轴定理是
\begin{equation}
 I_O=I_{\rm CM}+M(a^2\mathbf 1-\bm a\bm a\tp).
\end{equation}
其中$\mathbf 1$为单位矩阵。特别地，两条平行轴相距垂直距离$d$时，$I_O=I_{\rm CM}+Md^2$。
式\eqref{eq:cm-inertia}的纯转动形式也适用于惯性系中固定的支点，
但不能将任意加速点都当成固定支点使用所有转动公式。

\section{Euler方程与随体坐标}
选固定在刚体上的主轴。此坐标系中$I$不随时间变化，
但把角动量的惯性系导数换成随体导数时必须加上$\bm\omega\times\bm\ell$：
\begin{equation}
 \bm\tau=I\dot{\bm\omega}+\bm\omega\times(I\bm\omega).
\end{equation}
分量形式为
\begin{align}
 I_1\dot\omega_1+(I_3-I_2)\omega_2\omega_3&=\tau_1,\\
 I_2\dot\omega_2+(I_1-I_3)\omega_3\omega_1&=\tau_2,\\
 I_3\dot\omega_3+(I_2-I_1)\omega_1\omega_2&=\tau_3.
 \label{eq:cm-euler-rigid}
\end{align}
这里的力矩和惯量均关于质心，或关于一个惯性系中的固定支点。
点表示随体系分量对时间的导数。把$\bm\tau=I\dot{\bm\omega}$用于一般三维刚体会漏掉基矢转动。

\section{自由刚体与中间轴不稳定}
无外力矩时，惯性系中的$\bm\ell$固定，转动能$T$也守恒：
\begin{equation}
 \ell^2=I_1^2\omega_1^2+I_2^2\omega_2^2+I_3^2\omega_3^2,
 \qquad 2T=I_1\omega_1^2+I_2\omega_2^2+I_3\omega_3^2.
\end{equation}
在随体角速度空间中，运动沿这两个二次曲面的交线进行。
注意“角动量固定”说的是空间方向，随体分量仍会变化。

设$0<I_1<I_2<I_3$，考察绕第3主轴的定常转动$\omega_3=\Omega\ne0$。
保留横向扰动的一阶项，
\begin{equation}
 \dot\omega_1=\frac{I_2-I_3}{I_1}\Omega\omega_2,\qquad
 \dot\omega_2=\frac{I_3-I_1}{I_2}\Omega\omega_1,
\end{equation}
因而
\begin{equation}
 \ddot\omega_1=-\Omega^2
 \frac{(I_3-I_2)(I_3-I_1)}{I_1I_2}\omega_1.
\end{equation}
扰动振荡，所以绕最大惯量主轴的转动线性稳定；绕最小惯量主轴同样稳定。
绕中间轴时，
\begin{equation}
 \ddot\omega_1=\Omega^2
 \frac{(I_2-I_1)(I_3-I_2)}{I_1I_3}\omega_1,
\end{equation}
出现指数增长。这是网球拍翻转现象的局部机制，并非外力矩突然改变了空间角动量。
以上讨论为无耗散理想刚体。若有内部耗散但总角动量近似不变，
系统倾向于降低转动能，因此两种无耗散稳定旋转不一定具有相同的长期命运。

对自由对称陀螺$I_1=I_2=I_\perp$，$\omega_3$恒定，且
\begin{equation}
 \dv{}{t}(\omega_1+\i\omega_2)
 =\i\frac{I_3-I_\perp}{I_\perp}\omega_3(\omega_1+\i\omega_2).
\end{equation}
这是角速度横向分量相对随体轴的转动频率。
对称轴在空间绕固定角动量方向的进动频率则为$\ell/I_\perp$，两者不能混同。

\section{重对称陀螺：循环坐标降低难度}
考虑一端固定的轴对称刚体，质心距支点$l$，对称轴取由支点指向质心的方向。
以竖直向上的轴为参考，$\theta$为对称轴与竖直轴的夹角，$\phi$为进动角，$\psi$为自转Euler角。
惯量取关于支点的$I_1=I_2=I_\perp$及$I_3$，则
\begin{equation}
 L=\frac12I_\perp(\dot\theta^2+\dot\phi^2\sin^2\theta)
 +\frac12I_3(\dot\psi+\dot\phi\cos\theta)^2-Mgl\cos\theta.
\end{equation}
由于$\phi,\psi$循环，
\begin{equation}
 p_\psi=I_3(\dot\psi+\dot\phi\cos\theta)=\ell_3,\qquad
 p_\phi=I_\perp\dot\phi\sin^2\theta+\ell_3\cos\theta=\ell_z
\end{equation}
守恒。$\ell_z$是空间竖直角动量，$\ell_3$是随体对称轴分量。
消去两个角速度后，能量变成一维章动问题：
\begin{equation}
 E=\frac12I_\perp\dot\theta^2+U_{\rm eff}(\theta),\quad
 U_{\rm eff}=\frac{(\ell_z-\ell_3\cos\theta)^2}{2I_\perp\sin^2\theta}
 +\frac{\ell_3^2}{2I_3}+Mgl\cos\theta.
\end{equation}
对于$0<\theta<\pi$，有效势的转向点给出章动范围；Euler角在两极退化，需单独取极限分析。
此处由能量守恒降维是合法的，但把守恒动量直接代回$L$并对剩余坐标变分一般不合法；
固定动量的变分降维需要Routh变换。

匀速进动满足$\dot\theta=0$、$\dot\phi=\Omega_{\rm p}$。
在$\sin\theta\ne0$时，$\theta$方程给出
\begin{equation}
 I_\perp\Omega_{\rm p}^2\cos\theta-\ell_3\Omega_{\rm p}+Mgl=0.
\end{equation}
高速自转、慢进动支满足$|I_\perp\Omega_{\rm p}\cos\theta|\ll|\ell_3|$，于是
\begin{equation}
 \Omega_{\rm p}\simeq\frac{Mgl}{\ell_3}.
\end{equation}
这不是任意初态的精确解；通常还会叠加章动，快速进动支也被此近似排除。

\section{滚动：能量法与摩擦条件必须同时满足}
半径$R$、质心惯量$I=\beta mR^2$的轴对称圆柱状刚体，其质心在几何中心，
在倾角$\alpha$的固定斜面上无滑动滚动，转轴平行于斜面。
取沿斜面向下为正，静摩擦力大小$f$沿上坡方向，
\begin{equation}
 ma=mg\sin\alpha-f,\qquad I\dot\omega=fR,\qquad a=R\dot\omega.
\end{equation}
解得
\begin{equation}
 a=\frac{g\sin\alpha}{1+\beta},\qquad
 f=\frac{\beta}{1+\beta}mg\sin\alpha.
\end{equation}
无滑动成立还要求$f\leq\mu_smg\cos\alpha$。
在固定斜面、理想刚体且无滚阻时，接触点瞬时速度为零，静摩擦不做功，
若从静止下降高度$h$，能量法给出相同结果：$mgh=\tfrac12m(1+\beta)v^2$。
若斜面本身运动，或接触发生形变耗散，不能沿用这一功率判断。

\paragraph{复习与阅读}
必须说明惯量关于哪个点、分量在哪个系中、以及是否存在固定支点。
自由刚体优先用守恒量和Euler方程，重对称陀螺则用循环坐标降维。
延伸阅读：Tong讲义第3章；Taylor的刚体转动。
